% Copyright 2026 Alan Szmyt. All rights reserved pending publication license.
\section{Introduction}

Generative audio systems increasingly allow people to express musical intent
through natural language and other high-level controls. At the same time,
research in music psychology, affective computing, adaptive music, and adjacent
clinical contexts suggests that musical experiences can interact meaningfully
with subjective state. What remains unclear is whether a person's own
interpretable semantic vocabulary for desired internal transitions can be
mapped through generative audio into measurable acoustic realizations and then
refined longitudinally from observed response.

This manuscript treats that question as a research direction rather than an
established finding. The originating N-of-1 catharsis experience is classified
as an \textbf{observation} and a \textbf{hypothesis generator}; it is not a
result of the proposed study and does not establish efficacy, mechanism, or
generalizability.

\subsection{Central research question}

Can an adaptive system learn an individual mapping from interpretable sonic
language to generated acoustic structure to affective response, and use that
mapping to better target future state transitions?

\subsection{Background and novelty boundary}

The literature scan remains incomplete. Music and affect regulation,
personalized music systems, controllable generative audio, adaptive or
closed-loop systems, psychedelic and ketamine music research, and semantic
interpretability are the current review streams.

MindMelody is an important adjacent preprint because it describes a closed-loop
system that mediates EEG-derived affect through a semantic planning layer and
controllable music generation \citep{zhang2026mindmelody}. Its existence is
evidence that this design neighborhood is active prior art. Its reported pilot
results are not evidence that Antidote is effective, and the Antidote novelty
boundary must be tested against the full paper and additional primary sources.

\subsection{Candidate contribution}

\textbf{Hypothesis.} A useful contribution may be an interpretable,
user-extensible semantic control layer that records intended semantic controls,
the acoustic realization actually generated, subjective and optionally
physiological response, longitudinal relationships among those variables, and
uncertainty and provenance. This is a candidate contribution to be refined or
rejected after the novelty scan.

\BeaconFigure{semantic-acoustic-response-loop}{Candidate longitudinal mapping among semantic intent, realized audio, and individual response.}{Three boxes form a loop. Semantic intent points to generated acoustic structure, which points to measured individual response. A feedback arrow returns from response to the next semantic and generation decision. Each connection is labeled as a hypothesis to be tested.}

\section{Methods}

The candidate first study is a methods, system, and N-of-1 feasibility study,
not a clinical efficacy study. The protocol, variables, measures, controlled
mutations, analysis plan, and safety boundaries will be frozen only after the
primary-source literature and novelty scan.

\subsection{Conceptual conditional model}

Let $A$ denote generated audio, $S$ the current state, $T$ the desired state or
trajectory, $H$ the individual's prior response history, and $C$ the semantic
and creative control context. The working model considers
\begin{equation}
  P(A \mid S, T, H, C).
  \label{eq:conditional-audio}
\end{equation}
After exposure at trial $t$, audio $A_t$ is associated with measured response
$R_t$. The history may then be updated with the tuple
$(S_t, T_t, C_t, A_t, R_t)$. Equation~\ref{eq:conditional-audio} is a conceptual
model, not an identified statistical model or demonstrated mechanism.

\subsection{Planned measurement layers}

The minimum record should preserve semantic controls, model identity and
version, generation settings, output identity and checksum, acoustic features,
baseline and target states, subjective response measures, timestamps, protocol
version, and mutations between trials. Behavioral and physiological measures
are optional extensions; they must not become the sole definition of a person's
experienced state.

\subsection{Evidence classification}

Research records use five non-interchangeable classes: \textbf{source} for what
external evidence establishes, \textbf{hypothesis} for what the project proposes,
\textbf{observation} for what occurred, \textbf{interpretation} for a bounded
explanation of an observation, and \textbf{claim} for a statement supported in
the manuscript. The working claim ledger is maintained outside publication
artifacts under \texttt{research/notes/}.

\section{Results}

No formal study has been run and no formal results have been collected. The
originating catharsis experience is retained in the research bootstrap as an
observation, not reported here as a study result. This section will remain
explicitly empty until evidence is collected under a frozen protocol.

\section{Discussion}

The proposed framing emphasizes the longitudinal relationship between semantic
intent, measured acoustic realization, and an individual's response. That focus
may distinguish the work from systems centered primarily on population-level
affect labels, direct physiological conditioning, or one-time generation. This
is currently an interpretation of the design space, not a verified novelty
claim.

The term \emph{probability-space sculpting} is used as a computational metaphor:
semantic prompting and other controls constrain a conditional distribution over
possible generations rather than deterministically selecting one sound. The
metaphor becomes scientifically useful only if the intended controls and their
realized acoustic consequences can be operationalized and tested.

\section{Limitations}

The present concept does not establish treatment efficacy, universal mappings
between audio and mental states, or specific neurological mechanisms. Complex
mental-health conditions must not be reduced to a single brain region or audio
parameter. Binaural beats and related phenomena should be treated as testable
stimulus primitives with uncertain effects. Traditional or contemplative
ontologies may be studied as human meaning systems; they must not be presented
as established neuroanatomy.

The proposed N-of-1 design will have limited generalizability and substantial
risk of expectancy, demand, novelty, and measurement effects. Personalization
may overfit a short response history. Generative systems and acoustic feature
extractors may change over time. These risks require explicit versioning,
negative and null result retention, and conservative interpretation.

\section{Ethics statement}

No formal participant study or human-subject data collection is reported in
this draft. Any future protocol must resolve informed consent, privacy, data
retention, adverse-event handling, researcher-participant role conflicts, and
the applicable independent ethics-review requirements before collection.
Antidote is not a substitute for professional care and must not present
generated audio as treatment.

\section{Data and code availability}

No formal study data or prototype code accompanies this draft. The manuscript,
bibliography, migration inventory, and non-sensitive research records are
versioned in the Antidote repository. Generated publication artifacts are
reproducible through the Beacon dependency pinned in
\texttt{dependencies/beacon.lock.toml}. Human-subject or sensitive personal data
must not be committed to the public repository.

\section{Acknowledgements}

No external funding or institutional support is claimed in this draft.

\section{Contributor statement}

Contributor roles are generated from the project metadata using the CRediT
taxonomy. The current draft assigns conceptualization, methodology, original
draft writing, and review and editing to Alan Szmyt.
